%% kapa: exam=1 round=37 year=2026 period=1 form=A subject=real_estate subject_ko=부동산학원론 seq=19 number=99
%% source: Q-Net 감정평가사 기출문제 (cst00309) artlSeq=5258924 file=raw/1008/5258924/2026년 제37회 감정평가사 1차 1교시.pdf
%% conversion: kordoc 4.24.51 -> tools/gen1.py (auto); review=proofread-sonnet; flags=none
\begin{question}{99}
부동산시장과 부동산정책에 관한 설명으로 옳은 것은 모두 몇 개인가?
\begin{framedtext}
○ 용도지역ㆍ지구제와 토지비축제는 정부가 간접적으로 시장에 개입하는 제도이다. \par
○ 지가공시제는 실거래가 신고제보다 법령상 도입순서가 빠르다. \par
○ 주거급여제와 토지임대부 분양주택제는 현재 시행되고 있다. \par
○ 부동산시장에서 나타나는 부정적 외부효과는 시장실패의 요인이지만, 긍정적 외부효과는 시장실패의 요인이 아니다. \par
○ 재건축부담금은 종료시점 주택가액에서 개시시점 주택가액과 개발비용을 공제한 금액에 부과율을 적용하여 산정한다. \par
\end{framedtext}
\begin{choices}
  \item 없음
  \item 1개
  \item 2개
  \item 3개
  \item 4개
\end{choices}
\end{question}
\answer{3}
